\documentclass[12pt]{article}
\usepackage{graphicx,url}
\usepackage[brazil]{babel} 
\usepackage{lipsum}
%\usepackage[latin1]{inputenc}  
\usepackage[utf8]{inputenc}  %permite texto com acento
% UTF-8 encoding is recommended by ShareLaTex
\usepackage{verbatim}
\usepackage{listings}
\usepackage{xcolor}
\usepackage[margin=2.5cm]{geometry}
\usepackage{xcolor}
\usepackage{fancyhdr}
\pagestyle{fancy}
\fancyhf{}
\lhead{\it{\textcolor{gray}{\small V Mostra Científica de Tecnologia - UNIGRANDE}}}
\rhead{\it{\textcolor{gray}{\small Fortaleza-CE, 23 e 24 de março de 2021}}}
\lfoot{\it{XIII Feira Tecnológica - UNIGRANDE}}
\renewcommand{\headrulewidth}{0.5pt}

%%%%%%%%%%%%%%%%%%%%%%%%%%%%

\begin{document}
%\onehalfspacing

%
%%% Não alterar o preâmbulo acima.

    %%% MARQUE UM X NO TIPO DE PESQUISA


oindent{TIPO DE TRABALHO: 1. Estudo de caso( ), 2. Pesquisa bibliométrica(X), 3. Trabalho original( ), 4. Pesquisa em andamento do Trabalho de Conclusão de Curso( ).\\TIPO DE APRESENTAÇÃO: 1. AMOSTRA CIENTÍFICA(X) 2. SESSÃO TÉCNICA( ).
}
%%%%%%%%%%% TÍTULO %%%%%%%%%%%

\begin{center}
\textbf{\Large{TÍTULO DO TRABALHO}}\\
\end{center}

\vspace*{0.2cm}
%%%%%%%%%%% AUTORES - NO MÁXIMO 4 E 1 ORIENTADOR - %%%%%%%%%%%

\begin{flushright}
 {\bf Nome Completo do Aluno} \footnote[1]{Graduando em Ciência da Computação - FGF. e-mail: \it aluno@aluno.fgf.edu.br}  \\
 {\bf Nome Completo do Aluno} \footnote[1]{Graduando em Ciência da Computação - FGF. e-mail: \it aluno@aluno.fgf.edu.br} \\
  {\bf Nome Completo do Aluno} \footnote[2]{Graduando em Ciência da Computação - FGF. e-mail: \it aluno@aluno.fgf.edu.br}  \\
   {\bf Maria Edineuda Teixeira P. Nascimento} \footnote[3]{Mestre - Centro Universitário da Grande Fortaleza. e-mail: \it edineuda@fgf.edu.br}   \\
\end{flushright}

\vspace*{0.5cm}

%%%%%%%%%%% CORPO DO TRABALHO - DE 3 A 6 PÁGINAS %%%%%%%%%%%
%%%% TODOS OS TRABALHOS DEVEM TRAZER AS SEÇÕES EXATAMENTE COMO DESTACADO ABAIXO %%%%%%



oindent{\textbf{RESUMO:} Deve conter uma síntese sobre o que trata o trabalho, quais são os objetivos deste estudo e qual a natureza da pesquisa. Além disso, deve descrever apenas o título dos trabalhos relacionados que dão base à pesquisa. Por fim, deve informar brevemente qual é o método utilizado, os resultados (parciais ou finais) e, em síntese, qual foi a conclusão.

\vspace*{0.5cm}


oindent{\textbf{PALAVRAS CHAVES:} Deve conter até 3 palavras chaves.} 

%%%%%%%%%%% FIM DO RESUMO %%%%%%%%%%%
%%%%%%%%%%% O RESUMO DEVE FICAR APENAS NA PRIMEIRA PÁGINA %%%%%%%%%%%


\section{Introdução}


%\lipsum[2] % OBS.: Este comando gera texto aleatório. Deve apagá-lo para começar a escrever seu texto.
Este trabalho foi desenvolvido


\section{Trabalhos relacionados} \label{sec:firstpage}

Segundo \cite{1} ...


\section{Fundamentação teórica}

\lipsum[2] % OBS.: Este comando gera texto aleatório. Deve apagá-lo para começar a escrever seu texto.


\section{Procedimentos metodológicos}

\lipsum[2] % OBS.: Este comando gera texto aleatório. Deve apagá-lo para começar a escrever seu texto.


\section{Resultados}\label{sec:analisedosdados}

\lipsum[2] % OBS.: Este comando gera texto aleatório. Deve apagá-lo para começar a escrever seu texto.


\section{Conclusões}

\lipsum[2] % OBS.: Este comando gera texto aleatório. Deve apagá-lo para começar a escrever seu texto.


\begin{thebibliography}{1}

\bibitem{1}{IEEEhowto:kopka}
H.~Kopka and P.~W. Daly, \emph{A Guide to \LaTeX}, 3rd~ed.\hskip 1em plus 0.5em minus 0.4em\relax Harlow, England: Addison-Wesley, 1999.

\bibitem{2}{@inproceedings{inproceedings,
author = {Castro, Ana and Bento, Anderson and Rodrigues, João and Fonseca, Lucas and Silva, Alisson and Oliveira, Tiago and Lopes, Adriano and Pereira, Jean},
year = {2016},
month = {10},
pages = {},
title = {DESENVOLVIMENTO DE ROBÔ INTELIGENTE PARA A OLIMPÍADA BRASILEIRA DE ROBÓTICA}
}}


\end{thebibliography}


\end{document}


